% ১৩. প্রোগ্রাম ডিজাইন মডেল
\section{প্রোগ্রাম ডিজাইন মডেল}

\begin{learningbox}
এই টপিক শেষে শিক্ষার্থী structured, visual, object-oriented ও event-driven programming model-এর ধারণা তুলনা করতে পারবে।
\end{learningbox}

\begin{center}
\begin{tabular}{|p{0.24\textwidth}|p{0.42\textwidth}|p{0.22\textwidth}|}
\hline
\textbf{Model} & \textbf{ধারণা} & \textbf{উদাহরণ} \\
\hline
Structured & sequence, selection, loop; module-based & C \\
\hline
Visual & GUI control drag-and-drop করে design & Visual Basic \\
\hline
OOP & class ও object ভিত্তিক & C++, Java \\
\hline
Event-driven & user/system event অনুযায়ী কাজ & GUI app \\
\hline
\end{tabular}
\end{center}

\begin{center}
\fbox{\parbox{0.88\textwidth}{
\centering
\textbf{Model map placeholder}\\[4pt]
Problem type $\rightarrow$ Suitable programming model $\rightarrow$ Program design
}}
\end{center}

\begin{boardbox}{বোর্ড ফোকাস}
Programming model-এর প্রশ্নে শুধু নাম লিখবে না; কোন ধরনের problem/application-এ ব্যবহার হয় সেটিও লিখবে।
\end{boardbox}
